\chapter*{Introduction générale}
\addcontentsline{toc}{chapter}{Introduction générale}
\markboth{Introduction générale}{Introduction générale}
\label{chap:introduction}

À l'ère de la transformation numérique, le commerce électronique connaît une mutation profonde, portée par l'intégration croissante de l'intelligence artificielle. Les consommateurs ne se contentent plus de simples catalogues en ligne ; ils exigent des expériences personnalisées, immersives et intelligentes. Cependant, sur le marché local tunisien, un fossé subsiste entre ces attentes et les solutions proposées par les plateformes existantes, souvent limitées techniquement et fonctionnellement.

C’est dans cette optique que nous avons conçu \textbf{Alpha Store}, une plateforme e-commerce de nouvelle génération. Ce projet ambitionne de redéfinir l'expérience d'achat en ligne en intégrant des modules d'intelligence artificielle avancés — allant des algorithmes de recherche classiques comme le BFS et le DFS aux modèles de langage modernes tels que Gemini — au sein d'un écosystème multi-secteurs couvrant la mode et la technologie.

L’objectif principal de ce projet est de démontrer comment l'IA peut non seulement optimiser les processus de vente, mais surtout enrichir le parcours utilisateur à travers des fonctionnalités innovantes telles que l'essayage virtuel, l'optimisation budgétaire par algorithmes génétiques ou encore l'assistance conversationnelle intelligente.

Ce rapport détaille les différentes étapes de conception et de réalisation de ce projet, structurées en quatre chapitres :

\begin{itemize}[label=$\bullet$]
    \item \textbf{Le chapitre 1} présente le contexte général du projet, la problématique, l'étude de l'existant ainsi que la planification adoptée.
    \item \textbf{Le chapitre 2} est dédié à l'analyse des besoins et à la spécification fonctionnelle, incluant les cas d'utilisation et les maquettes de l'interface.
    \item \textbf{Le chapitre 3} expose la conception technique de la plateforme, avec les diagrammes de classes et l'architecture de la base de données.
    \item \textbf{Le chapitre 4} détaille la réalisation technique, les technologies employées (PHP MVC, Flask, GSAP), l'implémentation des modules d'IA et les tests de validation.
\end{itemize}

Enfin, nous conclurons ce rapport en synthétisant les résultats obtenus et en proposant des perspectives d'évolution pour Alpha Store.
